\documentclass{ximera}
\title{Potentiële elektrische energie, potentiaal en spanning}
\author{Daan Lipkens}

\begin{document}
\begin{abstract}
    Inleiding tot de potentiële elektrische energie van een lading in een elektrisch veld, de potentiaal in een punt en de spanning tussen twee punten.
\end{abstract}
\maketitle
\label{hoofdstuk:potentiaal}
% ----------------------------------------------------------------------
%          inleiding hoofdstuk
% ----------------------------------------------------------------------
\section*{Inleiding}
In de mechanica leerde je dat een massa in een zwaarteveld potentiële energie bezit, en dat deze energie kan worden omgezet in kinetische energie wanneer de massa valt. Iets gelijkaardigs geldt voor een lading in een elektrisch veld: ook hier is er sprake van potentiële energie, die kan worden omgezet in kinetische energie wanneer de lading zich onder invloed van de elektrische kracht verplaatst.

In dit hoofdstuk beperken we ons tot het \emph{homogeen} elektrisch veld tussen twee evenwijdige platen, zoals ingevoerd in het vorige hoofdstuk. We onderzoeken eerst hoe je de potentiële elektrische energie van een lading in zo'n veld kan berekenen. Vervolgens voeren we de \emph{potentiaal} in: een grootheid die, in tegenstelling tot de potentiële energie, niet afhangt van de grootte van de lading die je in het veld plaatst, maar enkel van het veld zelf en de plaats in dat veld.

Daarna bekijken we het verschil in potentiaal tussen twee punten, de \emph{spanning}. Spanning is een grootheid die je uit het dagelijks leven kent (denk aan de spanning van een batterij of van het stopcontact), en die nauw verbonden is met de arbeid die een elektrisch veld op een lading verricht. We sluiten af met een toepassing: de beroemde proef van Millikan, waarmee de elementaire lading experimenteel werd bepaald.

% \begin{definition}[title={Herhaling: elektrostatica}]\nl
% Dit hoofdstuk behoort, net als de vorige hoofdstukken over elektrische krachtwerking en het elektrisch veld, tot de \emph{elektrostatica}: het gedeelte van de fysica dat ladingen in rust bestudeert.
% \end{definition}

\end{document}
